\section{模型检验与评价}

\subsection{模型检验}

四问结果均通过统一的一致性检查：每个自然日含144个10 min时段和145个SOC状态；功率进入能量平衡前乘$1/6$；充电增加库存时乘$\eta_c$，放电减少库存时除以$\eta_d$；SOC与充放电功率不越界；跨日SOC按原浮点值连续传递；正常费用、调整费用与5倍紧急费用均由执行轨迹独立重算一致；扰动未来实际数据或尚未发布预报不改变当前决策。问题三另查提前量、发布时间、CVaR离散分位质量和可信度反演；问题四另查对数价格残差、依赖收缩、条件分组、桥接库存、LP净化和真实价格账单。

\subsection{模型优点}

本文用同一组物理约束贯穿四问，变量含义和量纲一致；按发布时间构建信息集，区分可执行策略与完美信息下界；联合残差日块保留可解释的时序轨迹和跨变量依赖；主体为线性或混合整数线性结构，便于解释、复现和审计。第二问通过事前验证选择预测器，第三问把提前量校准、概率场景、预报更新和尾部风险拆成可单独检验的模块，第四问则把价量依赖估计与下一节点前瞻设置端点和退化对照，使增量价值来源可追溯。

\subsection{局限与改进}

问题二固定0:00的名义电池轨迹，未利用日内已实现误差反馈；较强线性终端价值使333个正式日日末SOC贴近上限；随机森林也没有天气和节假日等外生变量。问题三90\%区间覆盖率仅80.67\%，且条件分箱尺度的平均费未优于统一尺度，说明有限日块仍需更长样本检验。问题四中依赖收缩和价格修正的费用差额区间均跨0，$K=3$虽在正式期较优却未通过事前验证，说明一年样本不足以支持精细价量结构的稳定选择。后续可采用48 h滚动时域、非线性库存价值、外生天气与日历变量，并用跨年度或保持季节性的块自助法检验。

\subsection{主要结论}

问题一验证储能在已知日型下的跨时段套利能力；问题二表明，因果预测、联合残差场景和储能计划能显著减少高倍紧急购电，但收益不可拆分归因给单一模块；问题三表明6:00与12:00的新预报具有主要增量价值，CVaR可以用较高平均费用换取较低尾部紧急成本。问题四表明场景方案相对点预测显著减少高价缺电，但价量依赖收缩和日内价格修正的细小增益尚无稳定统计证据；允许日内更新的冻结主方案比仅日前主方案少1.341\%，该差额应解释为预报信息和调整权限的合成价值。
